\chapter{Discusión}
\label{cap:discusion}

\section{Interpretación de los resultados}
\label{sec:disc_interpretacion}

Los resultados obtenidos permiten aceptar las cuatro hipótesis formuladas
en el apartado \emph{Objetivos} (\hyperref[sec:hipotesis]{Hipótesis de
partida}) con matices. La tasa de éxito del LLM en
la curación clínica ha sido del 85,9\,\% global
(figura~\ref{fig:curacion}), con 9 de 11 cohortes por encima del 95\,\%,
superando ampliamente el umbral del 80\,\% planteado en la hipótesis~1.
La firma génica validada por consenso multi-cohorte alcanza AUC 0,968 en
LODO sobre la tarea de subtipo (tabla~\ref{tab:lodo_sub}), muy por encima
del 0,90 planteado en la hipótesis~2. La recuperación del panel IHC es del
90\,\% (18/20), superando el 75\,\% planteado en la hipótesis~3. Y el
panel mínimo alcanza AUC 0,966 con 20 genes, a menos de 0,01 de la firma
completa (figura~\ref{fig:panel_minimo}), cumpliendo la hipótesis~4.

Más allá del cumplimiento formal de las hipótesis, los resultados aportan
tres observaciones cualitativas relevantes:

\begin{enumerate}
  \item \textbf{La recuperación del panel IHC escamoso es completa
  (12/12).} El framework, seleccionando genes por un criterio puramente
  estadístico (consenso de $d$ de Cohen en tres cohortes independientes),
  encuentra los mismos genes que la patología clínica utiliza en
  inmunohistoquímica diagnóstica desde hace décadas. Esto es una
  triangulación externa fuerte: la firma no es sólo estadísticamente
  robusta, es biológicamente correcta.
  \item \textbf{La discrepancia entre AUC y balanced accuracy en tumor
  \emph{vs.} sano} (visible en la figura~\ref{fig:auc_vs_balacc} y en la
  tabla~\ref{tab:lodo_tvs}) pone de manifiesto una limitación de la firma
  tal cual se entrega, pero también un problema técnico bien conocido y
  corregible. El modelo ordena bien las muestras (AUC 0,925) pero necesita
  recalibración del umbral entre cohortes. Este comportamiento es
  característico del desbalance del entrenamiento y no una limitación
  intrínseca del framework.
  \item \textbf{La existencia de dos ejes biológicos identificables}
  (pérdida alvéolo-capilar y proliferación, sección~\ref{sec:res_biologia})
  enriquece la interpretación clínica: la firma no es una caja negra,
  sino que corresponde a mecanismos oncológicos conocidos. La correlación
  inversa entre score y contenido de pulmón sano residual dentro de los
  tumores (figura~\ref{fig:composicion}) es una prueba adicional de que
  la firma captura señal biológica y no un confundente de composición.
\end{enumerate}

\section{Contribución del modelo de lenguaje}
\label{sec:disc_llm}

La integración de Llama~3.3-70b en el paso de curación clínica ha
demostrado ser una decisión de diseño acertada. Sin este componente, la
integración de 11 cohortes con nomenclaturas heterogéneas habría requerido
codificar reglas de mapeo específicas para cada estudio (o incluso para
cada muestra) o realizar la curación manualmente. Ambos caminos son poco
escalables y frágiles frente a la incorporación de nuevas cohortes.

Es importante subrayar que el LLM cumple aquí una función bien
delimitada: leer texto libre y devolver una etiqueta canónica.
\emph{No participa en ninguna decisión analítica posterior}. No selecciona
genes, no computa métricas, no interpreta resultados. Esta restricción es
deliberada y responde a dos preocupaciones:

\begin{itemize}
  \item \textbf{Sesgo por conocimiento paramétrico.} Un LLM entrenado en
  la literatura biomédica \emph{sabe} qué genes se han asociado
  clásicamente al cáncer de pulmón. Si se le permitiera participar en la
  selección de biomarcadores, sería difícil discernir si un gen aparece
  en la firma porque los datos lo señalan o porque el modelo lo
  recomienda.
  \item \textbf{Alucinación.} Los LLMs pueden inventar cifras plausibles
  cuando no tienen datos suficientes. En una tarea de razonamiento
  abierto sobre resultados de un análisis, este riesgo es real. La tarea
  de clasificación de texto (categoría de 3 opciones) es mucho menos
  susceptible de alucinación que la generación libre.
\end{itemize}

La misma restricción se aplica al asistente conversacional integrado en
la interfaz: su prompt sistema le prohíbe inventar cifras y le exige
responder \emph{«Eso no figura en los resultados del trabajo»} cuando no
tenga el dato. En el capítulo~\ref{cap:resultados} se muestra que el
asistente y el resto de la interfaz emiten cifras coherentes entre sí,
gracias a que ambos se alimentan del mismo conjunto de CSV y JSON del
pipeline.

\section{Sesgos y limitaciones metodológicas}
\label{sec:disc_sesgos}

\subsection{Ámbito de plataforma}
El clasificador se entrena sobre datos de microarray, principalmente
plataforma Affymetrix HG-U133 Plus 2.0. Su transferibilidad directa a
datos de RNA-Seq requiere normalización específica y no está garantizada:
las distribuciones de intensidad y el ruido técnico son distintos. Este
punto se discute como línea de trabajo futuro en la
sección~\ref{sec:conc_futuro}.

\subsection{Ámbito histológico}
El entrenamiento se restringe a adenocarcinoma y carcinoma escamoso
(tabla~\ref{tab:lodo_sub}). El framework no está diseñado para clasificar
tumores neuroendocrinos ni otras histologías raras
(figura~\ref{fig:histologias}). Esta restricción se declara explícitamente
en la interfaz y en el asistente conversacional, para evitar cualquier uso
inadecuado en triage sobre casos sin diagnosticar.

\subsection{Umbral de decisión}
Como se ha señalado, el umbral óptimo del clasificador varía entre
cohortes. Para uso clínico esto exigiría una calibración por cohorte de
implementación o una recalibración isotónica sobre un pequeño conjunto
de referencia local. Ambas soluciones son estándar en el despliegue de
clasificadores biomédicos, pero no se han integrado en este trabajo.

\subsection{Población y sesgos demográficos}
Las cohortes GEO empleadas se generaron mayoritariamente en poblaciones de
origen europeo o asiático. La transferibilidad a otras poblaciones no
está evaluada. Este es un sesgo común en genómica pública que la
comunidad reconoce~\cite{sirugo2019missing}.

\subsection{Tamaño muestral}
Dos cohortes evaluables (GSE23066 y GSE7670) tienen tamaños muy reducidos
($n \leq 15$, tabla~\ref{tab:cohortes}). Como se observa en la
tabla~\ref{tab:lodo_tvs}, precisamente estas cohortes son las que
concentran las AUC individuales más bajas (0,440 y 1,000 respectivamente),
reflejando la baja precisión estadística que impone su tamaño. El diseño
podría beneficiarse de restringir la evaluación a cohortes con al menos
30-50 muestras, aunque a coste de reducir el número de folds LODO.

\section{Comparación con la literatura}
\label{sec:disc_literatura}

Los resultados obtenidos se sitúan en el rango alto de la literatura
publicada sobre clasificación de subtipos NSCLC a partir de expresión
génica. Girard et al.~\cite{girard2016comprehensive} publican una
clasificación de subtipos moleculares en NSCLC con AUC en torno a 0,92
sobre TCGA, sin validación externa en cohortes independientes. Chen et
al.~\cite{chen2017gene} proponen una firma pronóstica de 12 genes con AUC
0,83 en validación externa. La firma presentada en este trabajo obtiene
AUC 0,968 en validación externa sobre tres cohortes independientes, con un
panel mínimo de 20 genes. La tabla~\ref{tab:comparativa} sitúa la firma
frente a algunas de las más relevantes publicadas en las últimas dos
décadas.

\begin{table}[H]
  \centering
  \small
  \caption{Comparación con firmas génicas publicadas para cáncer de
  pulmón.}
  \label{tab:comparativa}
  \begin{tabular}{lllrrl}
    \toprule
    \textbf{Firma} & \textbf{Año} & \textbf{Tarea} & \textbf{$n$ genes} & \textbf{AUC} & \textbf{Validación} \\
    \midrule
    Beer et al.~\cite{beer2002gene}           & 2002 & Pronóstico ADC           & 50   & --    & Interna (1 cohorte) \\
    Chen et al.~\cite{chen2017gene}           & 2007 & Pronóstico NSCLC estadio I & 5    & 0,83  & Externa (1 cohorte) \\
    Kratz et al.~\cite{kratz2012practical}    & 2012 & Pronóstico NSCLC no escamoso & 14  & --    & Externa (2 cohortes) \\
    Girard et al.~\cite{girard2016comprehensive} & 2016 & Subtipo NSCLC           & \emph{n.\,d.} & 0,92 & Interna (TCGA) \\
    Faruki et al.~\cite{faruki2017lung}       & 2017 & Subtipo NSCLC            & \emph{n.\,d.} & --    & Múltiples cohortes \\
    \midrule
    \textbf{Este trabajo}                     & 2026 & \textbf{Subtipo ADC vs SQC} & \textbf{1174} & \textbf{0,968} & \textbf{LODO estricto (3 cohortes)} \\
    \textbf{Este trabajo (panel mínimo)}      & 2026 & \textbf{Subtipo ADC vs SQC} & \textbf{20}   & \textbf{0,966} & \textbf{LODO estricto (3 cohortes)} \\
    \bottomrule
  \end{tabular}

  \vspace{0.35em}
  {\footnotesize \emph{n.\,d.}: no disponible en el artículo original.
  ``Validación externa'' se refiere a cohortes no usadas en la fase de
  descubrimiento; ``LODO estricto'' añade la garantía de que ninguna
  muestra de la cohorte test aparece en el entrenamiento.}
\end{table}

Más importante que la magnitud del AUC es el criterio de validación. Las
firmas publicadas suelen mostrar caídas de 5-15 puntos porcentuales entre
validación interna y externa~\cite{bernau2014crossvalidation}. La firma
presentada aquí ya está validada externamente, por lo que no cabe
esperar una degradación adicional al aplicarla a nuevas cohortes GPL570.

Respecto a la validación IHC, el trabajo de Faruki et
al.~\cite{faruki2017lung} ya señalaba la correspondencia entre firmas
transcriptómicas y marcadores IHC clásicos, pero sin cuantificarla ni
proponer un panel mínimo derivado. La aportación específica de este TFM
es cuantificar la recuperación (18/20) y explicitar el subconjunto no
recuperado (SFTPA1, SFTPC) con una explicación biológica plausible.

\section{Implicaciones clínicas potenciales}
\label{sec:disc_clinico}

El panel mínimo de 20 genes tiene, en principio, aplicabilidad clínica
directa como complemento diagnóstico en casos histológicamente ambiguos.
Su tamaño lo hace compatible con PCR cuantitativa multiplex (kits de hasta
30 tránscritos son rutinarios en laboratorios de patología molecular) o
con \emph{NanoString nCounter} (que trabaja habitualmente con paneles de
100-800 genes). El tiempo de respuesta ronda las 24-48\,h desde biopsia,
compatible con la práctica oncológica.

No obstante, dos requisitos deben cumplirse antes de considerar un uso
clínico real:

\begin{enumerate}
  \item \textbf{Validación prospectiva} sobre una cohorte no incluida en
  el descubrimiento, idealmente de una plataforma distinta (RNA-Seq) para
  demostrar transferibilidad tecnológica.
  \item \textbf{Aprobación regulatoria} (marcado CE-IVD en la UE, 510(k)
  o PMA en EE.UU.) que exigirá estudios clínicos formales y un
  fabricante que asuma la responsabilidad del ensayo.
\end{enumerate}

Ambos pasos exceden el alcance de este TFM, pero el panel identificado es
un candidato razonable para iniciar dicha ruta.

\section{Comparación con firmas comerciales existentes}
\label{sec:disc_comerciales}

Las firmas comerciales existentes para NSCLC (Pervenio Lung RS, Oncotype
DX Lung) se orientan principalmente al pronóstico de recurrencia tras
cirugía, no al diagnóstico histológico. La firma presentada en este
trabajo cubre un nicho diferente: la clasificación del subtipo
histológico. Por tanto, no compite directamente con las firmas
pronósticas, sino que las complementa cubriendo una necesidad clínica
menos atendida.

\section{Lecciones metodológicas}
\label{sec:disc_lecciones}

Del trabajo se derivan varias lecciones metodológicas aplicables a otros
proyectos de integración transcriptómica multi-cohorte:

\begin{itemize}
  \item \textbf{LODO como estándar mínimo.} Los resultados en validación
  interna sobre-estiman consistentemente el rendimiento externo. LODO
  debería ser la validación de referencia para cualquier firma que
  aspire a uso clínico.
  \item \textbf{Consenso multi-cohorte como filtro de robustez.}
  Exigir que un gen replique en varias cohortes independientes elimina
  la mayoría de falsos positivos derivados de sesgos técnicos
  específicos de una cohorte.
  \item \textbf{No corregir efecto lote conjuntamente cuando LODO se va
  a aplicar.} La corrección de lote enmascara la variabilidad que LODO
  precisamente pretende medir.
  \item \textbf{Delimitar el rol del LLM.} Un LLM bien acotado a una
  tarea de curación de texto es una herramienta escalable y fiable.
  Fuera de esa tarea (por ejemplo, en la selección de genes o en la
  redacción de conclusiones), sus riesgos superan sus beneficios.
  \item \textbf{Auditoría automática de coherencia.} La interfaz web
  incluye una capa de contexto que se lee de los CSV/JSON en cada
  render, lo que garantiza que las cifras mostradas siempre reflejen el
  estado actual del análisis. Este patrón, aplicado también al prompt
  del asistente, elimina una fuente típica de errores en trabajos de
  este tipo (cifras hardcodeadas que envejecen mal).
\end{itemize}
